\chapter{Introduzione a Java}
\label{chap:introduzione-java}

Il \cref{chap:strutture-dati-adt} ha distinto le operazioni offerte
su un insieme di dati dai dettagli della loro implementazione.
Questa separazione tra interfaccia e realizzazione ricompare nella
programmazione a oggetti, che costituisce una caratteristica
fondamentale di Java. Dopo aver studiato la programmazione con Python
nel \cref{chap:introduzione-python}, si esaminano ora il processo di
sviluppo in Java, la struttura di un programma e l'uso iniziale di
classi e oggetti.\footnote{Materiale di riferimento: S. Canazza,
\emph{Java}, slide del corso di Fondamenti di Informatica,
\nolinkurl{2026_FI_07_canazza_Java.pdf}.}

\section{Portabilità, compilazione ed esecuzione}
\label{sec:fi-java-esecuzione}

Java fu sviluppato a partire dal 1991 presso Sun Microsystems,
con un gruppo di progettisti guidato da James Gosling. Tra gli
obiettivi del linguaggio vi è la possibilità di eseguire uno stesso
programma su calcolatori con architetture differenti. Questa
\emph{portabilità} riguarda non soltanto il codice sorgente,
ma anche una rappresentazione compilata indipendente dalla CPU.

Il processo usuale comprende tre fasi: scrittura del \emph{codice
sorgente} in un file con estensione \texttt{.java}, compilazione
in \emph{bytecode} ed esecuzione mediante la \emph{Java Virtual
Machine} (JVM). Il \emph{Java Development Kit} (JDK) fornisce
fra gli altri gli strumenti necessari a compilare ed eseguire
il programma, in particolare i comandi \texttt{javac} e \texttt{java}.

\begin{definition}[Bytecode e macchina virtuale Java]
\label{def:fi-java-jvm}
Il \emph{bytecode Java} è un insieme di istruzioni destinate alla
JVM, non direttamente alla CPU fisica. La \emph{JVM} è l'ambiente
software che carica, verifica ed esegue tali istruzioni, interpretandole
o traducendole in codice macchina appropriato alla piattaforma.
\end{definition}

Compilando \nolinkurl{HelloTester.java} con \texttt{javac}
si ottiene normalmente \nolinkurl{HelloTester.class}.
L'esecuzione si avvia con \texttt{java HelloTester}:
si indica il nome della classe, non quello del file \texttt{.class}.

\begin{unipdterminal}
$ javac HelloTester.java
$ java HelloTester
Hello, World!
\end{unipdterminal}

Il compilatore traduce il sorgente una volta, mentre la JVM
esegue il bytecode sul calcolatore di destinazione e può caricare
anche classi della libreria standard. Una JVM può avvalersi di un
compilatore \emph{just-in-time} (JIT), che traduce durante
l'esecuzione parti del bytecode in istruzioni native. Lo stesso
file \texttt{.class} può quindi essere eseguito su piattaforme
compatibili dotate di una JVM appropriata, senza ricompilarlo
per ciascuna CPU. Questa è l'idea espressa dalla formula
\emph{compile once, execute everywhere}.

L'esecuzione del bytecode può combinare interpretazione e compilazione
JIT, con costi e possibilità di ottimizzazione differenti. Non esiste
tuttavia un ordine di efficienza valido per ogni programma e
ogni piattaforma.

La distinzione tra compilazione e interpretazione, introdotta
nella \cref{sec:fi-linguaggi-esecuzione}, non identifica due
categorie necessariamente esclusive: Java combina entrambe le fasi.
Anche CPython compila normalmente il sorgente in bytecode,
ma la fase è generalmente gestita in modo automatico
(\cref{sec:fi-python-esecuzione}). Un eseguibile nativo contiene
invece istruzioni per una specifica architettura; un file
\texttt{.class} contiene istruzioni per la JVM. La portabilità
non significa che ogni programma sia indipendente dall'ambiente:
restano necessari una JVM compatibile e le eventuali risorse
richieste dal programma.

\section{Il primo programma: classe e metodo principale}
\label{sec:fi-java-primo-programma}

Un semplice programma Java visualizza un messaggio sullo
\emph{standard output}, ossia l'uscita testuale del processo.
Si consideri il contenuto del file \texttt{HelloTester.java}:

\begin{unipdcode}[language=Java]
public class HelloTester {
    public static void main(String[] args) {
        // Scrive il messaggio sull'uscita standard.
        System.out.println("Hello, World!");
    }
}
\end{unipdcode}


In Java il codice di un'applicazione è organizzato in \emph{classi}.
La dichiarazione \texttt{public class HelloTester} introduce una
classe accessibile da altre classi. La parola riservata
\texttt{class} indica la dichiarazione di una classe, mentre
\texttt{public} è uno \emph{specificatore di accesso}. Il nome
di una classe pubblica di primo livello deve coincidere esattamente
con quello del file sorgente che la contiene, senza l'estensione:
in un file \texttt{HelloTester.java} si dichiara quindi
\texttt{public class HelloTester}. Un file sorgente può contenere
al più una classe pubblica di primo livello.

Il corpo della classe è racchiuso tra parentesi graffe. Al suo
interno si trova il \emph{metodo} \texttt{main}, che nella
forma tradizionale adottata nel corso costituisce il punto di ingresso
dell'applicazione e descrive le istruzioni da eseguire all'avvio.
La sua intestazione è
\begin{center}
  \texttt{public static void main(String[] args)}.
\end{center}
Qui \texttt{public} rende il metodo accessibile;
\texttt{static} indica che il metodo appartiene alla classe e
può essere invocato senza creare una sua istanza;
\texttt{void} indica che non restituisce alcun valore;
\texttt{String[] args} è un parametro che può ricevere
come array di stringhe gli argomenti forniti all'avvio del
programma. Nelle applicazioni introdotte qui il parametro
può rimanere inutilizzato.

Le parole riservate del linguaggio, tra cui \texttt{public},
\texttt{class}, \texttt{static} e \texttt{void}, hanno una funzione
sintattica stabilita e non possono essere usate come identificatori.
Java distingue lettere maiuscole e minuscole: \texttt{HelloTester},
\texttt{hellotester} e \texttt{HELLOTESTER} sono nomi diversi.

\subsection{Istruzioni, commenti e formato del sorgente}
\label{sec:fi-java-sintassi}

Il corpo di un metodo contiene \emph{istruzioni}, generalmente
terminate da un punto e virgola. Le parentesi graffe delimitano
blocchi di istruzioni e le parentesi tonde racchiudono gli argomenti
delle invocazioni di metodo. A differenza dell'indentazione di
Python (\cref{sec:fi-python-condizioni}), in Java rientri e
interruzioni di riga non delimitano i blocchi; sono comunque
essenziali per rendere leggibile il programma. Gli spazi rimangono
necessari quando separano elementi lessicali distinti.

Un commento introdotto da \texttt{//} prosegue fino alla fine
della riga. I delimitatori \texttt{/*} e \texttt{*/} racchiudono
invece un commento su una o più righe. I commenti non sono
istruzioni eseguite dal programma e possono documentare
la ragione di una scelta implementativa.

\clearpage

\section{Uscita standard e invocazione dei metodi}
\label{sec:fi-java-metodi-uscita}

Nell'istruzione
\begin{center}
  \lstinline|System.out.println("Hello, World!");|
\end{center}
\texttt{System} è una classe della libreria standard di Java.
Il suo campo statico \texttt{out} contiene un riferimento a un
oggetto della classe \texttt{PrintStream}, che gestisce l'uscita
standard. Il metodo \texttt{println} di tale oggetto scrive
il valore ricevuto e termina la riga; \texttt{print} scrive
senza aggiungere un ritorno a capo.

\begin{unipdcode}[language=Java]
System.out.print("Risultato: ");
System.out.println(3 + 4);
\end{unipdcode}

\begin{unipdterminal}
Risultato: 7
\end{unipdterminal}

Le istruzioni dell'esempio sono da collocare all'interno di un
metodo, come \texttt{main}. Il valore \texttt{"Risultato: "}
è una \emph{stringa letterale}, cioè una sequenza di caratteri
delimitata da virgolette doppie. Senza virgolette, il testo verrebbe
interpretato secondo le regole sintattiche del linguaggio.
\texttt{println} può ricevere anche valori numerici o il risultato
di un'espressione.

\begin{definition}[Invocazione di un metodo]
\label{def:fi-java-invocazione}
L'\emph{invocazione} di un metodo di istanza richiede un oggetto
ricevente, il nome del metodo e un elenco di argomenti tra
parentesi tonde. Nell'espressione
\texttt{oggetto.nomeMetodo(argomento1, argomento2)} gli argomenti
sono separati da virgole; se non ve ne sono, le parentesi restano
obbligatorie.
\end{definition}

Il nome del metodo e i tipi degli argomenti devono essere
compatibili con la sua dichiarazione; non è sufficiente conoscere
soltanto il nome del metodo. Le classi già disponibili e i loro
metodi formano la \emph{Java Platform API}: una libreria standard
che offre funzionalità utilizzabili senza conoscere o riscrivere
la loro implementazione. Questa separazione tra interfaccia e
realizzazione richiama l'astrazione discussa nella
\cref{sec:fi-adt-strutture}.

\section{Errori e ciclo di sviluppo}
\label{sec:fi-java-errori}

La compilazione verifica che il sorgente rispetti le regole del
linguaggio e che i nomi e le operazioni impiegati possano essere
risolti correttamente. Un errore lessicale o sintattico, come
una stringa priva delle virgolette di chiusura, impedisce di
ottenere il bytecode. Anche un nome di campo inesistente viene
segnalato dal compilatore: sostituendo \texttt{out} con
\texttt{aut} in \texttt{System.out.println(...)} si può ricevere
il messaggio \texttt{cannot find symbol}. In questo caso il testo
è sintatticamente riconoscibile, ma il simbolo richiesto non
esiste nella classe \texttt{System}; è quindi un errore
rilevato durante la compilazione, non un errore logico.

Anche la disponibilità delle classi e librerie necessarie
può dare luogo a errori nella fase di caricamento o
collegamento (\emph{linking}), distinta dalla verifica
sintattica del sorgente.

Un \emph{errore logico} si ha quando il programma è eseguibile
ma produce un risultato diverso da quello previsto. Ad esempio,
stampare \texttt{"Hell, World!"} al posto di
\texttt{"Hello, World!"} non viola alcuna regola del
linguaggio, ma non soddisfa la specifica del messaggio.

Un \emph{errore di esecuzione} (\emph{runtime error}) si
manifesta invece mentre il programma viene eseguito: la
lettura di un file non disponibile oppure la conversione
non valida di un testo inserito dall'utente possono, se non
gestite, interrompere l'esecuzione con un'eccezione.
Gli errori logici non provocano necessariamente un'eccezione,
e le eccezioni non indicano necessariamente un difetto
sintattico del programma. Un singolo errore nel sorgente può
anche produrre più messaggi del compilatore: è opportuno
esaminare anzitutto la prima causa segnalata.

Il lavoro di programmazione procede dunque attraverso il ciclo
\emph{modifica--compila--collauda}: dopo la modifica del sorgente,
si ricompila e si esegue il programma per verificarne il
comportamento. Il superamento della compilazione non dimostra
la correttezza del programma rispetto al problema da risolvere.

\section{Variabili, tipi e assegnamento}
\label{sec:fi-java-variabili}

Come in Python, una variabile permette di riutilizzare durante
l'esecuzione un valore. A differenza del modello di assegnamento
descritto nella \cref{sec:fi-python-variabili}, in Java una
variabile è dichiarata con un \emph{tipo}, che determina quali
valori possono esserle assegnati e quali operazioni sono ammesse.
La dichiarazione specifica il tipo e il nome; un'inizializzazione
può assegnare un valore contestualmente alla dichiarazione:

\begin{unipdcode}[language=Java]
int lit;
int monete = 15000;
double euro = 2.35;
String messaggio = "Valore totale in euro ";
\end{unipdcode}

L'\emph{assegnamento} modifica il valore di una variabile già
dichiarata: prima si valuta l'espressione a destra del segno
\texttt{=}, poi si memorizza il risultato nella variabile a
sinistra. Non è un'affermazione di uguaglianza matematica:
\texttt{totale = totale + euro;} aggiunge al totale il
valore corrente di \texttt{euro}. Per confrontare l'uguaglianza
fra valori primitivi si impiega invece \texttt{==}.

La dichiarazione di una variabile locale senza inizializzazione
non ne permette ancora la lettura: il compilatore richiede che
sia stata assegnata prima del suo utilizzo. La stessa variabile
locale non può essere dichiarata due volte nello stesso ambito;
può invece ricevere nuovi valori attraverso più assegnamenti.

\subsection{Identificatori e convenzioni}
\label{sec:fi-java-identificatori}

Un \emph{identificatore} è il nome attribuito a una variabile,
un metodo o una classe. In Java può contenere lettere e cifre,
oltre ad alcuni caratteri particolari, tra cui \texttt{\_} e
\texttt{\$}, ma non può iniziare con una cifra né coincidere
con una parola riservata o con i letterali
\texttt{true}, \texttt{false} e \texttt{null}. Un nome che
inizia con una cifra non potrebbe essere distinto, secondo le
regole lessicali del linguaggio, da un letterale numerico. Spazi,
separatori come \texttt{( ) \{ \} [ ] ; , .} e operatori
come \texttt{+}, \texttt{*}, \texttt{=}, \texttt{==} e
\texttt{\&\&} non possono fare parte di un identificatore.
Il carattere \texttt{\_} usato da solo non è ammesso come
identificatore nelle versioni moderne di Java.

Nella prassi si preferiscono nomi descrittivi che iniziano con
una lettera. Variabili e metodi utilizzano generalmente
l'iniziale minuscola, le classi l'iniziale maiuscola;
nei nomi composti si applica il \emph{camelCase}, per esempio
\texttt{totalEuro} e \texttt{HelloTester}. Poiché Java distingue
maiuscole e minuscole, è sconsigliabile utilizzare nomi che
variano soltanto per la capitalizzazione.

\subsection{Interi e numeri in virgola mobile}
\label{sec:fi-java-tipi-numerici}

I tipi primitivi \texttt{int} e \texttt{double} rappresentano,
rispettivamente, numeri interi con segno a 32 bit e numeri in
virgola mobile a doppia precisione a 64 bit. Le rappresentazioni
sottostanti richiamano il complemento a due e il formato
IEEE 754 studiati nel \cref{chap:rappresentazione-informazioni}.
I due tipi esprimono anche una distinzione nel significato dei
dati: una quantità di lire, priva di unità frazionarie, può essere
rappresentata da \texttt{int}; una somma di euro contenente
centesimi, in un esempio introduttivo, da \texttt{double}.
La rappresentazione \texttt{int} occupa meno spazio di
\texttt{double}; il confronto dei tempi di elaborazione dipende
invece dall'hardware e dalle operazioni effettuate.

Per scrivere un numero decimale nel sorgente Java si usa il
punto, senza separatori delle migliaia: \texttt{1936.27} è un
letterale \texttt{double}, così come \texttt{1.93E3}, che
rappresenta \(1.93\cdot10^3\). Gli operatori \texttt{*} e
\texttt{/} indicano prodotto e divisione. Se entrambi gli
operandi sono interi, la divisione \texttt{/} produce il
quoziente intero troncato verso zero; se almeno uno è
\texttt{double}, si esegue una divisione in virgola mobile.
Nel programma seguente \texttt{1936.27} fa quindi sì che il
risultato della divisione sia in virgola mobile.

\begin{unipdcode}[language=Java]
public class Coins3 {
    public static void main(String[] args) {
        int lit = 15000;
        double euro = 2.35;
        double dollars = 3.05;

        double totalEuro = lit / 1936.27;
        totalEuro = totalEuro + euro;
        totalEuro = totalEuro + dollars * 0.72;

        System.out.print("Valore totale in euro ");
        System.out.println(totalEuro);
    }
}
\end{unipdcode}


Le variabili rendono esplicito il significato dei diversi
contributi e permettono di aggiornare il calcolo senza
riscrivere tutte le costanti. In questo esempio l'assegnamento
\texttt{totalEuro = totalEuro + euro;} usa il valore già
conservato in \texttt{totalEuro}; una successiva assegnazione
ne modifica nuovamente il contenuto.

\section{Classi, oggetti e interfacce pubbliche}
\label{sec:fi-java-classi-oggetti}

I tipi primitivi sono adatti a dati elementari, mentre dati
più complessi sono rappresentabili mediante \emph{oggetti}.
Una \emph{classe} descrive un tipo di oggetti, specificandone
le operazioni e i dettagli necessari a realizzarle. Un oggetto
è un'\emph{istanza} di una classe. In questa fase interessa
soprattutto utilizzare classi già definite, distinguendo i
metodi accessibili dall'esterno dalla loro implementazione.

L'insieme delle operazioni pubbliche di una classe ne costituisce
l'\emph{interfaccia pubblica}. I metodi \texttt{public} possono
essere invocati dall'esterno nei contesti ammessi dal linguaggio;
i metodi \texttt{private} sono accessibili soltanto all'interno
della classe che li dichiara. Per usare un metodo pubblico non
occorre conoscere le istruzioni che ne realizzano il comportamento:
è sufficiente conoscere i suoi parametri, il risultato e le
condizioni di utilizzo. L'analogia con un ADT riguarda la
separazione tra comportamento e implementazione, non
l'identità fra una classe e un tipo di dato astratto.

\subsection{Parametri, risultati e metodi di \texttt{String}}
\label{sec:fi-java-metodi-string}

Un metodo di istanza opera su un oggetto \emph{ricevente},
indicato prima del punto nell'invocazione. Nella terminologia
delle slide tale oggetto costituisce il \emph{parametro implicito};
i valori indicati tra parentesi sono invece gli \emph{argomenti
espliciti}. Per esempio, in \texttt{greeting.length()} il
ricevente è \texttt{greeting} e non vi sono argomenti espliciti,
mentre in \texttt{river.replace("issipp", "our")} il
ricevente è \texttt{river} e gli argomenti espliciti sono due.

Un metodo può restituire un valore di un tipo stabilito oppure
essere dichiarato \texttt{void}, se non restituisce alcun
valore. \texttt{String.length()} restituisce un \texttt{int},
\texttt{String.replace(...)} una nuova \texttt{String},
mentre \texttt{System.out.println(...)} restituisce
\texttt{void}. Non è quindi lecito assegnare a una variabile
il risultato di \texttt{println}.

Le operazioni precedenti si possono applicare a più stringhe
nella stessa sequenza di istruzioni:

\begin{unipdcode}[language=Java]
String greeting = "Hello, World!";
int n = greeting.length();
String river = "Mississippi";
String bigRiver = river.toUpperCase();
String changedRiver = river.replace("issipp", "our");
\end{unipdcode}
Il valore di \texttt{n} è \texttt{13}. La variabile
\texttt{bigRiver} contiene \texttt{"MISSISSIPPI"}, mentre
\texttt{changedRiver} contiene \texttt{"Missouri"}.
Il metodo \texttt{replace} restituisce una nuova stringa,
senza modificare quella originale: \texttt{river} continua
a riferirsi a \texttt{"Mississippi"}. Queste istruzioni
vanno inserite nel corpo di un metodo.

L'\emph{intestazione} di un metodo comprende specificatori
di accesso, tipo restituito, nome e parametri dichiarati.
Per esempio, si consideri l'intestazione
\begin{center}
  \texttt{public String replace(String target, String replacement)}
\end{center}
che dichiara un metodo pubblico con due parametri di tipo
\texttt{String} e risultato dello stesso tipo. Il nome e i tipi dei
parametri costituiscono la \emph{firma} nel senso usato dalle
regole del linguaggio Java; il tipo restituito e gli
specificatori di accesso appartengono invece all'intestazione
completa. La descrizione dei parametri, dei valori restituiti
e delle eventuali eccezioni si trova nella documentazione
ufficiale della \emph{Java Platform API}, disponibile anche
all'indirizzo
\url{https://docs.oracle.com/en/java/javase/21/docs/api/index.html}.

\section{Variabili oggetto e riferimenti}
\label{sec:fi-java-riferimenti}

Una variabile di tipo primitivo contiene un valore, mentre una
variabile di tipo oggetto contiene un \emph{riferimento} a un
oggetto. Una dichiarazione come \texttt{Rectangle box;}
introduce una variabile di riferimento, ma non crea un
rettangolo né consente di usare \texttt{box} prima di
un'assegnazione valida. L'espressione
\texttt{new Rectangle(5, 10, 20, 30)} invoca un costruttore,
crea un oggetto e restituisce un riferimento che può essere
assegnato alla variabile \texttt{box}.

L'esempio utilizza la classe \texttt{java.awt.Rectangle},
che rappresenta un rettangolo mediante coordinate, larghezza
e altezza. L'istruzione \texttt{import java.awt.Rectangle;}
ne permette l'uso con il nome semplice \texttt{Rectangle}.

\begin{unipdcode}[language=Java]
import java.awt.Rectangle;

public class RectangleTester {
    public static void main(String[] args) {
        Rectangle box = new Rectangle(5, 10, 20, 30);
        Rectangle alias = box;
        alias.translate(15, 25);
        System.out.println(box.getX());

        Rectangle copia = new Rectangle(box);
        copia.translate(10, 0);
        System.out.println(box.getX());
        System.out.println(copia.getX());
    }
}
\end{unipdcode}
\begin{unipdterminal}
20.0
20.0
30.0
\end{unipdterminal}


L'assegnamento \texttt{alias = box} copia soltanto il
riferimento: i due nomi individuano lo stesso oggetto.
Il metodo \texttt{translate} ne modifica quindi lo stato
attraverso entrambi i riferimenti, come schematizzato nella
\cref{fig:fi-java-alias-copia}. Viceversa, la costruzione di
\texttt{new Rectangle(box)} crea un distinto oggetto con
lo stesso stato iniziale; una successiva modifica di
\texttt{copia} non modifica \texttt{box}. Il metodo
\texttt{getX()} restituisce un valore di tipo
\texttt{double}, ragione per cui l'output numerico
contiene \texttt{.0}.

\begin{figure}[htbp]
  \centering
  \begin{tikzpicture}[unipd/diagram,x=1cm,y=1cm]
  \node[unipd/node,minimum width=1.25cm] (box) at (-3.7,0.7) {\texttt{box}};
  \node[unipd/node,minimum width=1.25cm] (alias) at (-3.7,-0.45) {\texttt{alias}};
  \node[unipd/node,minimum width=1.25cm] (copy) at (-3.7,-2.0) {\texttt{copia}};
  \node[unipd/node,minimum width=3cm,minimum height=0.85cm,align=center] (first) at (1.1,0.15)
    {\texttt{Rectangle}\\\(x=20,\ y=35\)};
  \node[unipd/node,minimum width=3cm,minimum height=0.85cm,align=center] (second) at (1.1,-2.0)
    {\texttt{Rectangle}\\\(x=30,\ y=35\)};
  \draw[unipd/directed edge] (box.east) -- ++(0.9,0) |- ($(first.west)+(0,0.18)$);
  \draw[unipd/directed edge] (alias.east) -- ++(0.9,0) |- ($(first.west)+(0,-0.18)$);
  \draw[unipd/directed edge] (copy.east) -- ++(0.9,0) |- (second.west);
\end{tikzpicture}

  \caption{Due riferimenti allo stesso oggetto e una copia indipendente.}
  \label{fig:fi-java-alias-copia}
\end{figure}

Questo comportamento è diverso dall'assegnamento tra variabili
primitive: dopo \texttt{int a = 13; int b = a;},
la modifica \texttt{b = 12;} lascia \texttt{a} uguale a
\texttt{13}. La differenza dipende da \emph{che cosa viene
copiato}: il valore numerico nel primo caso, il riferimento
all'oggetto nel secondo. Si tratta di un principio essenziale
per utilizzare correttamente le classi senza confondere
l'identità di un oggetto con quella delle variabili che vi
fanno riferimento.

Il collaudo di una classe mediante un programma che ne costruisce
oggetti e invoca i metodi è sviluppato nella
\cref{sec:fi-java-collaudo}, insieme ai tipi fondamentali e
all'ingresso standard di Java.
